% Multiple Integrals 100: four subrectangles with the lower-left sample points (from Lab 11).
\documentclass[border=4pt]{standalone}
\usepackage{tikz}
\begin{document}
\begin{tikzpicture}[scale=2.1]
  \draw[->] (-0.1,0) -- (1.3,0) node[right] {$x$};
  \draw[->] (0,-0.1) -- (0,1.3) node[above] {$y$};
  \draw (0,0) -- (1,0) -- (1,1) -- (0,1) -- cycle;
  \draw (0.5,0) -- (0.5,1);
  \draw (0,0.5) -- (1,0.5);
  \filldraw[red] (0,0)     circle (.8pt);
  \filldraw[red] (0.5,0)   circle (.8pt);
  \filldraw[red] (0,0.5)   circle (.8pt);
  \filldraw[red] (0.5,0.5) circle (.8pt);
  \node[below] at (0.5,0) {$\frac{1}{2}$};
  \node[below] at (1,0)   {$1$};
  \node[left]  at (0,0.5) {$\frac{1}{2}$};
  \node[left]  at (0,1)   {$1$};
\end{tikzpicture}
\end{document}
